% Compile twice: xelatex Lecture_8_Flow_Measurement_Math.tex
\documentclass[aspectratio=169,11pt]{beamer}
\usepackage{fontspec,polyglossia,amsmath,amssymb,booktabs}
\setdefaultlanguage{ukrainian}
\setmainfont{DejaVu Sans}
\setsansfont{DejaVu Sans}
\setmonofont{DejaVu Sans Mono}
\newfontfamily\cyrillicfont{DejaVu Sans}
\newfontfamily\cyrillicfontsf{DejaVu Sans}
\newfontfamily\cyrillicfonttt{DejaVu Sans Mono}
\definecolor{ink}{HTML}{183341}
\definecolor{accent}{HTML}{087C83}
\setbeamercolor{normal text}{fg=ink,bg=white}
\setbeamercolor{frametitle}{fg=accent}
\setbeamertemplate{navigation symbols}{}
\setbeamertemplate{footline}[frame number]
\title{Лекція 8. Вимірювання витрати та вибір витратоміра}
\subtitle{Математичний конспект і практичні завдання}
\date{Новий план, редакція 3 від 24.10.2026}
\begin{document}
\begin{frame}\titlepage\end{frame}
\begin{frame}{Позначення та межі моделей}
\begin{itemize}
\item $Q_v$ --- об'ємна витрата, м$^3$/с; $Q_m$ --- масова витрата, кг/с.
\item $V$ --- об'єм, м$^3$; $m$ --- маса, кг; $t$ --- час, с.
\item $\rho$ --- густина, кг/м$^3$; $p,\Delta p$ --- тиск і перепад, Па.
\item $D,L,\ell$ --- довжини, м; $A,k$ --- площа та коефіцієнт, м$^2$.
\end{itemize}
Усі числові приклади, описи середовищ і паспорти задач синтетичні. Реальні інженерні приклади мають окремі джерела.

Чотири блоки: 15, 12, 13 і 12 хв; вступ із повторенням 4 хв; завершення 4 хв. Пошук відповідей поза цими 60 хвилинами.
\end{frame}
\section{Витрата, кількість і змінний перепад тиску}
\begin{frame}{Визначення: витрата та кількість}
\[
Q_v(t)=\frac{dV}{dt},\qquad Q_m(t)=\frac{dm}{dt}=\rho(t)Q_v(t).
\]
Остання рівність стосується густини й об'єму за тих самих умов у відповідному перерізі для однорідного середовища.
\[
V(t_b)-V(t_a)=\int_{t_a}^{t_b}Q_v(t)\,dt.
\]
Для навчальної кусково-сталої моделі $Q_v(t)=Q_i$ на $[t_i,t_{i+1})$:
\[
\Delta V=\sum_{i=0}^{n-2}Q_i(t_{i+1}-t_i).
\]
Це точна сума для прийнятої моделі й наближення для реального потоку. Останній показ не створює нового інтервалу без його кінця.
\end{frame}
\begin{frame}{Одиниці та опорні умови газу}
\[
1\;\text{л}=10^{-3}\;\text{м}^3,\qquad
1\;\frac{\text{л}}{\text{хв}}=\frac{10^{-3}}{60}\;\frac{\text{м}^3}{\text{с}}.
\]
\[
Q_m=\rho_{\rm роб}Q_{v,\rm роб}=\rho_{\rm оп}Q_{v,\rm оп}.
\]
Для ідеального газу незмінного складу:
\[
Q_{v,\rm оп}=Q_{v,\rm роб}
\frac{p_{\rm роб}}{p_{\rm оп}}\frac{T_{\rm оп}}{T_{\rm роб}}.
\]
Тиски $p$ абсолютні, у Па; температури $T$ у К; обидві витрати у м$^3$/с. Для реального газу потрібна поправка на стисливість.

\texttt{slm} позначає стандартні літри за хвилину: конкретні опорні умови визначає документація. Надлишковий тиск і градуси Цельсія не підставляють у ці відношення.
\end{frame}
\begin{frame}{Визначення: основа перепадного методу}
Горизонтальний стаціонарний потік нестисливої рідини, ідеалізація без втрат:
\[
p_1+\frac{\rho v_1^2}{2}=p_2+\frac{\rho v_2^2}{2},\qquad
A_1v_1=A_2v_2=Q_v.
\]
Тут $v_1,v_2$ --- середні швидкості, м/с; $A_1,A_2$ --- площі, м$^2$.
\[
Q_v=C_d\frac{A_2}{\sqrt{1-\beta^4}}\sqrt{\frac{2\Delta p}{\rho}}
=k\sqrt{\frac{\Delta p}{\rho}},\quad \beta=\frac{d}{D}.
\]
$d,D$ --- діаметри звуження й труби, м; $C_d,\beta$ безрозмірні, $0<\beta<1$; $k=C_d A_2\sqrt{2/(1-\beta^4)}$ має одиницю м$^2$.

$C_d$ враховує реальну характеристику для визначених умов. Сталий $k$ --- навчальне припущення. Для газів потрібні додаткові поправки, зокрема на розширення.
\end{frame}
\begin{frame}{Розмірність, крайні точки й обернення}
\[
\left[\frac{\Delta p}{\rho}\right]
=\frac{\mathrm{kg\,m^{-1}s^{-2}}}{\mathrm{kg\,m^{-3}}}
=\mathrm{m^2s^{-2}},\qquad [Q_v]=\mathrm{m^3s^{-1}}.
\]
За $k=0{,}001$ м$^2$, $\rho=1000$ кг/м$^3$:
\[
Q_v(0\;\text{Па})=0,\quad
Q_v(1000\;\text{Па})=0{,}001\;\text{м}^3/\text{с},
\]
\[
Q_v(4000\;\text{Па})=0{,}002\;\text{м}^3/\text{с}.
\]
Обернене перетворення для $Q_v\ge0$, $k>0$, $\rho>0$:
\[
\Delta p=\rho(Q_v/k)^2.
\]
Підстановка прямої формули відновлює $\Delta p\ge0$. Від'ємний перепад у цій односторонній моделі відхиляється, а не замінюється модулем.
\end{frame}
\begin{frame}{Нелінійність і монтаж}
Два однаково тривалі режими: $\Delta p_1=0$ Па, $\Delta p_2=4000$ Па. За попередніх $k,\rho$:
\[
\overline{Q_v}=\tfrac12(0+2)=1\;\text{л/с},\qquad
Q_v(\overline{\Delta p})=\sqrt2\approx1{,}414\;\text{л/с}.
\]
Отже, перетворення середнього сигналу не замінює усереднення витрат.

За незмінних $k,\rho$: $Q_v/Q_{v,\max}=\sqrt{\Delta p/\Delta p_{\max}}$. При 10\% витрати сигнал перепаду становить 1\% повної шкали.

Перепад між відборами й безповоротна втрата тиску різні. Імпульсні лінії можуть додати гідростатичні складові. Геометрія відборів, стан ліній і робочий діапазон належать до умов моделі.
\end{frame}
\section{Сталий перепад, турбінні й теплові витратоміри}
\begin{frame}{Визначення: ротаметр}
Для нерухомого поплавка у вертикальному потоці вгору:
\[
mg=\rho gV_{\rm п}+F_{\rm потоку}.
\]
$m$ --- маса поплавка, кг; $g$ --- прискорення, м/с$^2$; $V_{\rm п}$ --- витіснений об'єм, м$^3$; сили у Н.

У спрощеній моделі $F_{\rm потоку}\approx\Delta p A_{\rm еф}$, де $A_{\rm еф}$ --- ефективна площа, м$^2$:
\[
\Delta p\approx\frac{mg-\rho gV_{\rm п}}{A_{\rm еф}}.
\]
Підняття поплавка збільшує кільцевий прохід у конічній трубці. За майже сталої ефективної ваги перепад приблизно сталий, а положення залежить від витрати.

Густина, в'язкість, тертя та орієнтація впливають на характеристику. Універсальної лінійної шкали висоти для всіх середовищ немає.
\end{frame}
\begin{frame}{Визначення: імпульсний метод}
\[
f=\frac{N}{\Delta t},\qquad Q_v=\frac{f}{K},\qquad V=\frac{N}{K}.
\]
$N$ --- число імпульсів; $f$ --- імп/с; $K$ --- імп/м$^3$; $\Delta t$ --- с. Результати: $Q_v$ у м$^3$/с, $V$ у м$^3$.
\[
\delta Q_{\rm крок}=\frac{1}{K\Delta t}.
\]
За $K=20\,000$ імп/м$^3$ крок для вікна 1 с дорівнює $5\cdot10^{-5}$ м$^3$/с, тобто 3 л/хв; для 10 с --- 0,3 л/хв.

Це крок підрахунку, а не повна похибка сенсора. Довше вікно усереднює короткі зміни. Зворотний зв'язок: $N=KV=KQ_v\Delta t$; модель передбачає сталий калібрувальний коефіцієнт у заданому режимі.
\end{frame}
\begin{frame}{Тепловий метод: що входить у калібрування}
Потік змінює відведення тепла від нагрітого елемента або асиметрію температурного поля. Зв'язок сигналу з витратою залежить від конструкції, складу газу та умов.
\[
Q_m=\rho_{\rm оп} Q_{v,\rm оп}.
\]
$Q_m$ --- кг/с; $\rho_{\rm оп}$ --- кг/м$^3$; $Q_{v,\rm оп}$ --- м$^3$/с за визначених опорних умов.

SFM3100 є документованим прикладом теплового каналу для повітря й кисню з температурною компенсацією. Позначення стандартної об'ємної витрати не звільняє від перевірки опорних умов.

Перенесення калібрування на інший газ без перевірки не випливає з фізики методу. Ця лекція не встановлює універсальної степеневої формули для всіх теплових сенсорів.
\end{frame}
\section{Електромагнітні й ультразвукові витратоміри}
\begin{frame}{Визначення: електромагнітне перетворення}
Для ідеалізованої геометрії:
\[
E\approx B\ell\bar v,\qquad Q_v=\bar v A,\qquad A=\frac{\pi D^2}{4}.
\]
$E$ --- В; $B$ --- Тл; $\ell,D$ --- м; $\bar v$ --- середня по перерізу швидкість, м/с; $A$ --- м$^2$.
\[
[B\ell\bar v]=\mathrm{T\,m^2/s}=\mathrm{V}.
\]
За $B=0{,}02$ Тл, $\ell=D=0{,}05$ м, $\bar v=2$ м/с:
\[
E\approx0{,}002\;\text{В},\qquad Q_v\approx0{,}003927\;\text{м}^3/\text{с}.
\]
Обернення: $\bar v=E/(B\ell)$ за $B\ell\ne0$. Реальний сигнал просторово зважений, тому потрібне калібрування. Звичайна модель вимагає заповненої труби, придатної провідності та контакту електродів.
\end{frame}
\begin{frame}{Визначення: різниця часів прольоту}
За рівномірної швидкості $v_{\rm ш}$ уздовж вимірюваного шляху:
\[
t_+=\frac{L}{c+v_{\rm ш}\cos\theta},\qquad
t_-=\frac{L}{c-v_{\rm ш}\cos\theta},
\]
\[
v_{\rm ш}=\frac{L}{2\cos\theta}
\left(\frac1{t_+}-\frac1{t_-}\right).
\]
$L$ --- м; $c,v_{\rm ш}$ --- м/с; $t_+,t_-$ --- с; $\theta$ --- кут до осі потоку. Потрібні $\cos\theta\ne0$ та $c>|v_{\rm ш}\cos\theta|$.

За $v_{\rm ш}=0$ обидва часи рівні $L/c$, обернена формула дає нуль. Для потоку вздовж додатного напряму $t_+<t_-$. Перехід від швидкості на промені до середньої по перерізу потребує моделі профілю.
\end{frame}
\begin{frame}{Доплерівський метод і межі акустичної моделі}
Для спрощеної геометрії зворотного розсіювання та $|v_s|\ll c$:
\[
\Delta f\approx\frac{2f_0v_s\cos\theta}{c}.
\]
$f_0,\Delta f$ --- Гц; $v_s$ --- швидкість розсіювачів, м/с; $c$ --- м/с; кут $\theta$ безрозмірний. Знак залежить від прийнятого напряму й визначення частотної різниці.

Потрібні розсіювачі та зв'язок їхньої швидкості з потоком. Частинки можуть осідати, а бульбашки рухатися відносно рідини.

Часовий метод потребує достатньої якості прийнятого сигналу. Для накладного приладу додатково враховують стінку, заломлення, акустичний контакт і геометрію. Універсального допустимого відсотка бульбашок для всіх моделей немає.
\end{frame}
\section{Промисловий вибір, монтаж і калібрування}
\begin{frame}{Визначення: чутливість до одного входу}
Для $Q_v=k\sqrt{\Delta p/\rho}$, додатних базових входів:
\[
\frac{Q'_v}{Q_v}=\frac{k'}k\sqrt{\frac{\Delta p'}{\Delta p}\frac{\rho}{\rho'}}.
\]
За незмінних інших входів:
\begin{align*}
\Delta p'=1{,}04\Delta p&:\quad Q'_v/Q_v=\sqrt{1{,}04}\approx1{,}019804,\\
\rho'=1{,}04\rho&:\quad Q'_v/Q_v=1/\sqrt{1{,}04}\approx0{,}980581,\\
k'=1{,}01k&:\quad Q'_v/Q_v=1{,}01.
\end{align*}
Зміни результату відповідно $+1{,}98\%$, $-1{,}94\%$, $+1\%$.

Це аналіз чутливості, не повний бюджет невизначеності. Для нульового базового перепаду відносна формула не визначена: застосовують абсолютну модель.
\end{frame}
\begin{frame}{Паспортна база та монтаж}
Навчальний приклад для $Q_v=10$ м$^3$/год, $Q_{\max}=100$ м$^3$/год:
\[
0{,}005Q_v=0{,}05\;\text{м}^3/\text{год},\qquad
0{,}005Q_{\max}=0{,}5\;\text{м}^3/\text{год}.
\]
Однакове число 0,5\% не задає однакової абсолютної межі. Перевіряють базу, додаткові члени та діапазон застосування.

Для заданого паспортом $L_{\rm до}\ge nD$: $D$ і $L_{\rm до}$ у м, $n$ безрозмірне. Наприклад, $D=0{,}05$ м та $n=10$ вимагають 0,5 м. Номінальний DN сам по собі не є точним внутрішнім діаметром.

Інтелектуальний канал потребує перевірки вихідної змінної. Якщо сигнал уже пропорційний витраті, повторне кореневе перетворення некоректне.
\end{frame}
\begin{frame}{Калібрування зважуванням}
\[
\overline{Q_m}=\frac{m_b-m_a}{t_b-t_a},\qquad
\overline{Q_v}=\frac{\overline{Q_m}}{\rho}
\]
за сталої відповідної густини; маси у кг, час у с, густина у кг/м$^3$. У динамічній моделі оцінюють $dm/dt$.

Установка NIST використовує приріст маси з поправкою на виштовхувальну дію повітря. Реальний аналіз також ураховує невизначеності часу, маси й інших входів.

Калібрування встановлює зв'язок опорних значень і показів із їхніми невизначеностями та використовує його для отримання результату з показу (JCGM VIM 2.39). Регулювання є окремою дією. Результат калібрування не скасовує монтажних вимог.
\end{frame}
\begin{frame}{Інженерні історії: джерела для самостійного читання}
\footnotesize
\begin{enumerate}
\item ASME, \textit{Holyoke Water Power System}: облік води фабрик і витратомір Гершеля.\\
\url{https://www.asme.org/about-asme/engineering-history/landmarks/129-holyoke-water-power-system}
\item Sensirion, \textit{SFM3100 Datasheet}: тепловий канал для вентиляційного обладнання.\\
\url{https://sensirion.com/products/catalog/SFM3100}
\item Emerson, \textit{Reliable Measurement Without Contact}, MS-00830-0200-0003, 2025: польовий кейс Endeavor, повідомлення виробника.\\
\href{https://www.emerson.com/is/content/emerson/en/measurement-instrumentation/marketing/products/flow-non-intrusive-ultrasonic/documents/doc-flexim-00830-0200-0003.pdf}{Повний кейс Emerson (PDF)}
\item NIST, \textit{Liquid Flow Standards}; докладний звіт NIST SP 250-98 (2021).\\
\url{https://doi.org/10.6028/NIST.SP.250-98}
\end{enumerate}
Перевірено 27.10.2026. Числові вправи не відтворюють параметри цих історичних або промислових установок.
\end{frame}
\begin{frame}{Технічні джерела та екзаменаційний зв'язок}
\scriptsize
\href{https://www.hr.endress.com/hr/support-overview/learning-center/flow-measuring-principle-emf}{Endress+Hauser: електромагнітний принцип};
\href{https://www.hr.endress.com/hr/support-overview/learning-center/flow-measuring-principle-ultrasonic-sensors}{ультразвукові принципи}.

\href{https://www.krohne.com/en-vn/products/flow-measurement/flowmeters/variable-area-flowmeters}{KROHNE: ротаметри та монтажні умови}.
\href{https://www.germiona.com.ua/files/-------350.pdf}{Метран-350: керівництво виробника, архівна копія}; приклад із програми, без тверджень про поточний асортимент.
\href{https://jcgm.bipm.org/vim/en/2.39.html}{JCGM VIM 2.39: калібрування}.

\medskip
57. Одиниці вимірювання.\\
58. Вимірювання витрати за перепадом тиску.\\
59. Витратоміри сталого перепаду тиску.\\
60. Турбінні, теплові, електромагнітні та ультразвукові рівноміри.\\
61. Інтелектуальні витратоміри ПГ „Метран”.\\
62. Промислові типи витратомірів.

\medskip
Питання 60 відтворено дослівно: «рівноміри» в офіційному переліку, контекст теми --- витратометрія. Покриття: блок 1 --- 57--58; блок 2 --- 59 і частина 60; блок 3 --- завершення 60; блок 4 --- 61--62.
\end{frame}
\begin{frame}{Практичне завдання 1. Перепад і накопичений об'єм}
\small
Синтетичні дані: $t=[0,10,25,40]$ с, $\Delta p=[0,1000,4000,1000]$ Па; $k=0{,}001$ м$^2$, $\rho=1000$ кг/м$^3$.
\begin{enumerate}
\item Створіть функцію перерахунку перепаду у витрату й функцію накопичення. На кожному $[t_i,t_{i+1})$ прийміть витрату рівною лівому показу.
\item Поверніть список $Q_i$ у м$^3$/с, три об'єми інтервалів і загальний об'єм у м$^3$ та літрах. Останній показ не продовжуйте за межі запису.
\item Відхиляйте різні довжини списків, менш ніж дві точки, нескінченні або відсутні числа, незростаючий час, $k\le0$, $\rho\le0$, $\Delta p<0$.
\item Перевірте нульовий і сталий потік, нерівні інтервали та зворотне перетворення $\Delta p=\rho(Q_v/k)^2$ із числовим допуском.
\item Поясніть, чому підстановка середнього перепаду може змінити облік. Здайте код, одиниці та короткий звіт про тести.
\end{enumerate}
Базовий Python, функції, списки й \texttt{math}; готового рішення немає.
\end{frame}
\begin{frame}{Практичне завдання 2. Два вікна імпульсного обліку}
\small
Синтетичний коефіцієнт $K=20\,000$ імп/м$^3$; A: $N=7$, $\Delta t=1$ с; B: $N=73$, $\Delta t=10$ с. Це два окремі записи, а не обов'язково вкладені вікна.
\begin{enumerate}
\item Для кожного запису поверніть частоту, об'єм і середню витрату за вікно, а також крок $\delta Q_{\rm крок}$.
\item Подайте об'єм у м$^3$ та л, витрату в м$^3$/с та л/хв. Явно зазначте одиницю K.
\item Перевірте $N=0$, $N=1$, неціле або від'ємне N, $K=0$ і $\Delta t=0$; некоректні записи не перетворюйте на нульові результати.
\item Перевірте оберненість $KV=N$ та $KQ_v\Delta t=N$ для допустимих даних із допуском округлення.
\item Порівняйте швидкість оновлення й крок двох вікон. Поясніть, чому нуль у короткому вікні не доводить відсутності повільного потоку й чому крок не дорівнює повній похибці.
\end{enumerate}
Результат: функція, два словники виходу, тести та пояснення без спеціальних бібліотек.
\end{frame}
\begin{frame}{Практичне завдання 3. Відбір принципів}
\scriptsize
Синтетичні середовища, труба всюди повна, акустичний контакт придатний; $\sigma$ у мкСм/см, $b$ --- об'ємна частка газу у \%:
\begin{center}
\begin{tabular}{lrrl}
\toprule Середовище & $\sigma$ & $b$ & Розсіювачі, що рухаються з рідиною\\\midrule
Вода & 500 & 0 & ні\\
Олива & 0,01 & 0 & ні\\
Суспензія & 800 & 2 & так\\
Невідоме & невідомо & 0 & невідомо\\\bottomrule
\end{tabular}
\end{center}
Навчальні правила: електромагнітний E вимагає $\sigma\ge5$; часовий U вимагає $b\le0{,}5$ і відсутності розсіювачів; доплерівський D вимагає наявності розсіювачів. Це умовні обмеження задачі, не універсальні характеристики технологій.
\begin{enumerate}
\item Створіть таблицю рішень для всіх 12 пар «середовище--принцип».
\item Поверніть статус і всі причини. Відома невідповідність означає відмову; без неї невідоме поле означає недостатність даних; виконання всіх вимог означає придатність за правилами.
\item Перевірте пороги $\sigma=5$, $b=0{,}5$, відсутнє поле та поєднання невідомого з явно невиконаною вимогою. Поясніть обмеження отриманого відбору.
\end{enumerate}
\end{frame}
\begin{frame}{Практичне завдання 4. Паспорти й монтаж}
\scriptsize
Умови: чиста вода, $\sigma=500$ мкСм/см, $Q=2$--$20$ м$^3$/год, $T=40\,^{\circ}$C, робочий тиск 6 бар, DN50; повна труба; доступно 5D до й 3D після; допустима втрата при 20 м$^3$/год не більше 10 кПа. У всіх кандидатів DN50, вода та потрібний вихід підтверджені.
\begin{center}
\begin{tabular}{lrrr}
\toprule Синтетичний паспорт & A: електромагн. & B: турбінний & C: часовий\\\midrule
Витрата, м$^3$/год & 1--25 & 3--30 & 0,5--25\\
Температура, $^{\circ}$C & 0--80 & 0--60 & 0--70\\
Макс. тиск, бар & 10 & 16 & 10\\
Прямі ділянки до/після & 5D/2D & 10D/5D & 10D/3D\\
Втрата при 20 м$^3$/год, кПа & 2 & 15 & 0 (накладний)\\\bottomrule
\end{tabular}
\end{center}
Для A мінімальна $\sigma=5$ мкСм/см; для C стінка й акустичний контакт підтверджені. Інші вимоги цієї навчальної задачі вважаємо виконаними.
\begin{enumerate}
\item Відберіть кандидатів за всіма вимогами та поверніть всі причини відмов. Не порівнюйте лише максимальні витрати.
\item Додайте тест на рівність межі та тест із відсутнім полем. Складіть монтажний чекліст і перелік налаштувань вихідної змінної.
\item Поясніть, які зміни монтажу могли б змінити відбір, не ігноруючи інших невиконаних вимог. Здайте код і обґрунтування.
\end{enumerate}
\end{frame}
\end{document}
